% Isolated Version 14 author draft; assembled independent review is required.
% This file does not modify the original polynomial or Version 13.
\section{A quartic presentation with 299 variables}\label{sec:v14}

We retain the arithmetic inputs, qualifications and source editions of
Sections~\ref{subsec:v12-inputs} and~\ref{sec:v13}. The polynomial of
the main construction is unchanged. The following are smaller
presentations of the same rational input projection.
\begin{theorem}\label{thm:v14-quartic}
There is a specified integer polynomial $P_{299}$ of total degree four
in 299 variables such that
\[
 X\in\mathbb Z\quad\Longleftrightarrow\quad
 \exists\mathbf v\in\mathbb Q^{298}:P_{299}(X,\mathbf v)=0.
\]
It has 6,630 nonzero monomials, content one, and largest absolute
coefficient $186624\cdot3^{2^{26}}$. Its coefficient artifact writes
$C=3^{2^{25}}$ as a fixed integer, not a quantified variable.
There is also a specified integer quartic $P_{314}$ with
\[
 X\in\mathbb Z\quad\Longleftrightarrow\quad
 \exists\mathbf w\in\mathbb Q^{313}:P_{314}(X,\mathbf w)=0.
\]
It has 6,672 nonzero monomials and content one. Every coefficient has
absolute value strictly less than $2^{6496}$; its largest coefficient
has exactly 6,496 bits. The complete ordinary decimal coefficient table
occupies 222,890 bytes in the specified ASCII format.
\end{theorem}
The ordinary coefficient bound here is different from the bound for
$P_{411}$. We prove the new bound for the new artifact. Some changes
below are inverse changes of private witnesses; others change an
arithmetic certificate or point domain. For the latter we prove
soundness and completeness separately.

\subsection{Conventions and retained interfaces}\label{subsec:v14-inputs}
Put $K=\mathbb Q(\rho)$, $\rho^2=2$, $d=3-\rho$, $D=d^2=11-6\rho$,
and $E:y^2=x^3-Dx$. Retain
$E(K)=\mathbb ZR_0+E[2](K)$ and $2E(K)=\mathbb ZP$, where
\[
 R_0=(1+\rho,2\rho),\quad P=(c_*,a_*b_*),\quad
 a_*=-1+7\rho/4,\quad b_*=39/8-9\rho/2,\quad c_*=a_*^2.
\]
Thus $b_*^2=c_*^2-D$. Write $f(Q)=1/x(Q)$ and
$z(Q)=-x(Q)/y(Q)$, with both functions extended regularly by zero
at $O$. The roots are $Q_i=[n_i]P$ with labels
$(n_1,n_2,n_3,n_4)=(n,a,t,b)$. Their nonzero odd integral labels
will be restored before height or logarithm arguments. The scalar
point is $Q_2-P=[s]P$, where $s=a-1$.

The retained inputs include the squareclass calculation, formal
expansions and logarithm estimates, simultaneous contact-prime
construction with its real-density condition, preparation-zero pole
criterion, root-height implication, Lucas argument and scalar-depth-24
argument. In particular the three positive rational square scales are
\[
 A_0=1,\quad A_1=C_0^2,\quad A_2=(1+C_1^2)^2,\qquad
 \prod_i x(Q_i)=C_0+C_1\rho .
\]
The same root-height conclusion is
\[
 B^2-6<7\mathcal H|N|^{19/k},\qquad
 B=\max_i|n_i|,\quad 1\leq\mathcal H<2^{2093966},
\]
when the $k$th powers of the three required contact radicals divide
the nonzero integer $N$. The independently restored scalar anchor
gives $B>(40\mathcal H)^{10}$, excluding all three radicals equal
to one without using H or auxiliary logarithms.
The retained preparation-zero equations exclude negative odd
valuations of a pole input at required large primes. Their
completeness interface uses the same small-prime squareclasses and
depths, square reduced denominator, and real interval $(0,17)$.
We verify the changed consumers of these inputs below.

\subsection{The fixed-difference chart for H}\label{subsec:v14-H}
First replace the old $x(H)$ pair by the equations
\begin{equation}\label{eq:v14-H-raw}
 c_*(f+F)+Dv^2+2b_*v=1,\qquad v=ft,\qquad F=vt.
\end{equation}
The common primary normalizer is the fixed affine expression
\begin{equation}\label{eq:v14-normalizer}
 \mathcal S=(b_*+c_*t+Dv)/a_*.
\end{equation}
Its substitution into every primary, logarithm and scalar comparison
row preserves quadratic degree. If $f=0$, the last two equations
force $v=F=0$, contradicting the first. Set $x=1/f$. They give
$v^2=fF$ and $a_*f\mathcal S-b_*f-(c_*+Df)v=0$.
Here is an unconditional certificate for the recovered curve.
Put $d_1=c_*+Df$, $B_1=b_*f+d_1v$,
$A_1=a_*f\mathcal S-B_1$, and
$K_1=c_*(f+F)+Dv^2+2b_*v-1$. Then
\[
 c_*f(f\mathcal S^2+Df^2-1)
 =fd_1K_1+c_*d_1(v^2-fF)
 +(b_*^2-c_*^2+D)f^2+A_1(a_*f\mathcal S+B_1).
\]
Thus $\mathcal S^2=x-Df$ and the old fixed-difference equations
hold. This direction divides only by $c_*f$. Conversely, the old
full chart gives $t=v/f$ and the same affine expression for
$\mathcal S$; its previously established exceptional-point exclusion
supplies the forward implication. All outside values are preserved.

Next use the already needed outputs $U=f(H)^2$, $V=f(H-P)^2$.
Store $W,v,t,U,V$, put $Q=(1-DW-2b_*v)/c_*$, and impose
\begin{equation}\label{eq:v14-H-squared}
 W=v^2,\qquad UV=W^2,\qquad Q^2=U+V+2W,\qquad t(U+W)=Qv.
\end{equation}
On $UQ\ne0$, set $f=(U+W)/Q$, $F=(V+W)/Q$. If
$E_1=UV-W^2$ and $E_2=Q^2-U-V-2W$, then
\begin{align*}
 (U+W)^2-UQ^2&=-UE_2-E_1,\\
 (V+W)^2-VQ^2&=-VE_2-E_1,\\
 (U+W)(V+W)-WQ^2&=-WE_2+E_1.
\end{align*}
These give $f^2=U$, $F^2=V$, $fF=W$, $f+F=Q$, $ft=v$,
and $F=vt$, restoring~\eqref{eq:v14-H-raw}.

The extra branches are classified before this reconstruction.
For $U=0$ one has $W=v=0$, $Q=1/c_*$, $V=1/c_*^2$.
For $Q=0$ one has $U=V=-v^2$ and
\[
 Dv^2+2b_*v-1=0,\qquad
 v=\frac{3+2\rho}{4}\quad\hbox{or}\quad
 v=\frac{-36+16\rho}{49}.
\]
These are units at every place above $p>76176$:
$c_*+b_*=12-8\rho$ has norm16 and
$c_*-b_*=(9+4\rho)/4$ has norm$49/16$.
Section~\ref{subsec:v14-groups} excludes these branches by their
unchanged $V_0$ unit leaf and the pre-H height/anchor contradiction.
No division by $UQ$ is used to prove its own legitimacy.

\subsection{Root parity and the shifted anchor}\label{subsec:v14-parity}
The root chart $c_i=a_i^2$, $b_i^2=c_i^2-D$ represents
$(c_i,a_ib_i)$. Its abscissa is nonzero; for the final presentation
this follows from the retained product-height inverse as proved in
Section~\ref{subsec:v14-reciprocal}.
The squareclass map has independent values
\[
 \alpha(R_0)=[1+\rho],\qquad \alpha(T_0)=[-1],\qquad
 \alpha(T_d)=[d].
\]
Indeed norms in a putative square relation force the exponents
of $1+\rho$ and $d$ even, and the real embeddings exclude the
remaining class $[-1]$. Thus the square abscissa already puts
each root in $2E(K)=\mathbb ZP$, with nonzero integral label,
before an odd-chart condition.

Keep the scalar reciprocal equations
\[
 (c_*c_2+D)g=a_*b_2-b_*a_2,\qquad
 g=g_0+g_1\rho,\qquad g_0\tau=1.
\]
The denominator is nonzero at both real embeddings. The subtraction
identity gives $f(Q_2-P)=g^2$. At equal abscissae, $Q_2=P$
contradicts the guard and $Q_2=-P$ gives the exact duplication
value for $-2P$. Hence $s=n_2-1\ne0$ is an integer before height
or logarithms. Write $F_s=2g_0g_1$.

Let $X_n$ be the $\rho$ coefficient of $x([n]P)$. For $n\ne0$,
\begin{equation}\label{eq:v14-X}
 v_3(X_n)=
 \begin{cases}0&n\text{ odd},\\-2-2v_3(n)&n\text{ even},\end{cases}
 \qquad v_2(X_n)=-1\quad(n\text{ odd}).
\end{equation}
At3 the reduction of $P$ is $(\rho,0)$ over $\mathbb F_9$,
which proves the odd case. For $n=2\cdot3^ku$, $u$ a3-unit,
the formal expansion gives
$z([n]P)/3^{k+1}\equiv u(2+\rho)\pmod3$.
Since $x(z)=z^{-2}+O(z^2)$ and
$(2+\rho)^{-2}=3/2-\rho\equiv2\rho\pmod3$, the scaled
abscissa has nonzero imaginary coefficient modulo3. This proves
the even case and $X_n\ne0$.
At2, $z(P)=(232+86\rho)/357$ has component valuations3 and1.
The integral odd multiplication series preserves these values:
its error has field valuation at least$9/2$ when $v(2)=1$.
For $z=A+B\rho$ with those valuations the imaginary coefficient
of $z^{-2}$ is $-2AB/(A^2-2B^2)^2$, of valuation$-1$.
The Laurent correction has larger valuation. This uses integral
formal coefficients, not good reduction at2.

The retained scalar-coefficient lemma gives $v_3(F_n)=0$ for odd
$n$, $2+2v_3(n)$ for even nonzero $n$, and dyadic integrality.
It is applied with both parities. Put
\begin{equation}\label{eq:v14-q}
 q=12(X_A^2+X_T^2)+108X_Z^2,\qquad T=F_s+Cq,
\end{equation}
and use the private anchor
\begin{equation}\label{eq:v14-anchor}
 u^2+6v_a^2-2w_a^2-4Cw_a-3T^2=0.
\end{equation}
The inherited quaternion criterion, including parameter zero,
supplies rational witnesses exactly when $T/C$ is integral at2
and3. Its3-integrality forces all three summands in $F_s/C+q$
to be3-integral. Their possible finite valuations lie in the
disjoint classes $2\mathbb Z$, $1+4\mathbb Z$, $3+4\mathbb Z$;
the binary sum uses anisotropy modulo3 and~\eqref{eq:v14-X}.
A negative minimum cannot cancel. Thus $X_A,X_T,X_Z$ are3-units,
so the three root labels are odd. Also $F_s/C$ is3-integral
and is dyadically integral by the scalar lemma. The old unshifted
anchor can therefore be restored with new private witnesses
before H or height. An odd $s$ would make $F_s$ a3-unit,
so $s$ is even and the S-root label $s+1$ is odd.

For an odd root, its half point $[n_i]R_0=(A,B)$ has
$A=(1+\rho)r^2$. Duplication gives
\[
 a_{\rm half}=(A^2+D)/(2B),\qquad
 b_{\rm half}=(B/A)^2/2-c_i,
\]
with $a_{\rm half}^2=c_i$ and the correct product ordinate.
After changing both root signs if necessary, the old half-chart
equations are
$A^2-2(c_i+b_{\rm half})A-D=0$ and
$\operatorname N(A)+\operatorname N(r)^2=0$.
In the final A,T,Z modules the sign freedom is $a_i\mapsto-a_i$,
with reconstructed $b_i\mapsto-b_i$. For S negate
$a_2,b_2,g_0,g_1,\tau$ together. These actions fix the actual
points, $F_s$, all detector values and the private unscaled anchor
witnesses. The scalar rows change at most by an overall sign.

Conversely the old odd roots make $q$ integral at2 and3 by
\eqref{eq:v14-X}, so the same quaternion criterion supplies the
shifted anchor. No prime or H choice changes. The exact private
unscaling used here multiplies the old anchor witnesses and its
parameter $B$ by $g_0$, with inverse multiplication by $\tau$.
It gives $T=g_0B=F_s+Cq$ and~\eqref{eq:v14-anchor}.
This exact step assumes no private anchor witness nonzero.

\subsection{Contact poles and their conditional domains}\label{subsec:v14-contact}
For each scale put
\[
 R_j=F_5(A_j\lambda),\quad E_jR_j=1,\qquad
 F_5(t)=t(t+1/2)(t+1)(t-1/2399)(t+32/407).
\]
The retained contact gates make $R_j$ affine in existing coordinates.
With $W_j=G_{j,0}G_{j,1}G_{j,2}$ retain only
\begin{equation}\label{eq:v14-pole}
 \operatorname N(W_j)=2\operatorname{Re}(W_j)z_j,\qquad
 B_j(R_j+E_j+2)=z_j^2.
\end{equation}
There is no quantified inverse of $\operatorname{Re}(W_j)$.
At a required contact prime, $\alpha=v_p(R_j)>0$ is odd
and $1+R_j$ is a unit. The groups below give $W_j\ne0$.
The norm row then forces its real component and $z_j$ nonzero:
$w_0^2-2w_1^2=0$ over $\mathbb Q$ forces both components zero.
Thus
\[
 T_j^{\rm tr}=\operatorname{Tr}(W_j^{-1})=1/z_j,\qquad
 B_j=\frac{R_j}{(T_j^{\rm tr})^2(1+R_j)^2}.
\]
These domains are needed only at required contact primes.
A detector with radical one has empty local obligations;
no global trace inverse is asserted for it. The case $R_j=-1$
likewise cannot occur at a required contact.

In completeness the ratio between $R_j$ and the earlier contact
rational function is $A_j^5(t-1)^6$, a rational square and a
unit at selected contacts. Hence $R_j$ has valuation one there
and even valuation outside the prescribed odd support.
Both denominators in $B_j$ change valuations by even integers.
This includes additional primes introduced by later witnesses.

\subsection{Common-scale groups and the successful place}\label{subsec:v14-groups}
Use the physical tests
\[
 H_0=f(H),\quad H_i=L_{i+1}-\overline{LP_{i+1}}\ (1\leq i\leq3),
 \quad H_4=\mathcal L(L_2,L_3,L_4),
\]
\[
 M_0=L_1-\overline{LP_1},\quad M_1=L_1-\beta(X),\quad
 P_i=f(U_i)^2,\quad Z_1=\overline{f(V_1)^4},\quad
 T_i=\overline{f(V_i)^8}.
\]
Here $\beta(X)=204(X^2+X)+2555$ is the retained comparison
polynomial. Conjugation is part of the physical tests.
The Lucas polynomial is
\[
 \mathcal L(A,T,Z)=T^2-A(A-2)TZ+(A-2)^2Z^2-1.
\]
The affine recentering $C_L=(A-2)Z$, $V_L=AT$ gives
$\mathcal L+1=T^2+C_L^2-C_LV_L$ without changing its value.

For a nonzero common rational $\gamma$, store
\[
 ZH_i=H_i^4/\gamma^{160},\quad ZM_i=M_i^4/\gamma^{40},\quad
 ZP_2=P_2^2/\gamma^{160},\quad
 ZP_3=P_3^2/\gamma^{80},\quad ZP_4=P_4^2/\gamma^{80},
\]
and $ZP_1=P_1^2$. These definitions are quadratic using the
retained square inputs and common powers. For each detector,
\begin{align}
 G_0&=1+RZH_1+R^2ZH_2+EU+E^4Z_1+\gamma^{16}E^8,\label{eq:v14-g0}\\
 G_1&=1+RZH_3+R^2ZH_4+R^3ZP_2+E^4ZP_4+EZP_1+E^8T_0,\label{eq:v14-g1}\\
 G_2&=1+RZM_0+R^2ZM_1+R^3ZP_3+E^8T_2+E^4T_3+ET_4.\label{eq:v14-g2}
\end{align}
Here $U=H_0^2$ is the actual squared-H output even before
H is reconstructed. The product, trace and pole use these
same final groups.

At a consumed place set $\alpha=v(R)>0$ odd and
$c_\gamma=v(\gamma)\in\mathbb Z$. For a reciprocal point
$h^2=f(1-Df^2)$, a nonzero $f$ has even valuation: the leading
valuation is $v(f)$ when positive and $3v(f)$ when negative.
Zero leaves are omitted. Thus the relevant powers have
periods4,8,16 before formality or logarithms. In the generic
squared-H branch its algebraic reconstruction gives the same
curve fact; its exceptional branches have the fixed signed
values already classified. Other high and middle tests are
literal fourth powers, independently of their later ratio meanings.

In term order, the valuation periods and coefficients of
$\alpha$ are
\[
 \begin{array}{c|l}
 G_0&(4,1),(4,2),(4,-1),(8,-4),(16,-8)\\
 G_1&(4,1),(4,2),(8,3),(8,-4),(8,-1),(16,-8)\\
 G_2&(4,1),(4,2),(8,3),(16,-8),(16,-4),(16,-1).
 \end{array}
\]
All gamma exponents are divisible by the relevant periods.
Within each row the classes modulo16 are disjoint and exclude
zero for every odd $\alpha$. Every group therefore has a unique
minimum valuation, is nonzero, and has valuation at most zero.
This covers arbitrary $c_\gamma$ and every zero-leaf pattern.

Let $\ell_p$ be the maximum valuation of $W$ over places above
$p$. Trace integrality gives $v_p(T^{\rm tr})\geq-\ell_p$.
The pole obstruction and~\eqref{eq:v14-pole} give
\[
 0\leq v_p(B)=\alpha-2v_p(T^{\rm tr})\leq\alpha+2\ell_p.
\]
Its odd parity makes $\ell_p>-\alpha/2$. At a place attaining
the maximum, every group, and then every individual term,
has valuation greater than $-\alpha/2$. This is one place for
all tests, also at an inert unramified prime.

The guard gives $c_\gamma>15\alpha/32$. The high and middle
terms give
\[
 4v(H_i)>160c_\gamma-(r+1/2)\alpha,\qquad
 4v(M_i)>40c_\gamma-(r+1/2)\alpha\quad(r=1\text{ or }2).
\]
For $\alpha=1$ the worst bounds are157.5 and37.5; for
$\alpha=3$ they are312.5 and72.5; for $\alpha\geq5$ they
exceed$72.5\alpha$ and$16.25\alpha$. Thus $v(H_i)\geq40$
and $v(M_i)\geq10$. The primary bounds are
\[
 \begin{split}
 4v(f(U_2))&>160c_\gamma-3.5\alpha,\\
 4v(f(U_3))&>80c_\gamma-3.5\alpha,\\
 4v(f(U_4))&>80c_\gamma+3.5\alpha.
 \end{split}
\]
Using even valuations, these give respectively40,20,22 for the
reciprocal depths. The other low terms force positive valuation
of every required nonzero primary or alternate reciprocal value.

On either exceptional squared-H branch $T_0$ is a unit at every
consumed place. Its $E^8T_0$ term has valuation$-8\alpha$,
contradicting the common strict bound. Either branch would
therefore make all radicals one, contrary to the pre-H
root-height/anchor bounds. The unchanged fixed-point exclusions
from Section~\ref{sec:v13} then restore the full H/V0 domain
before logarithms.
